\item \textbf{Problema de repaso.} La figura P1.28 muestra a los estudiantes que analizan la conducción térmica de la energía en bloques cilíndricos de hielo. Este proceso se describe por la ecuación
$$ \frac{Q}{\Delta t} = \frac{k \pi d^2 (T_h - T_c)}{4L} $$
Para control experimental, en estos ensayos todas las cantidades, excepto $d$ y $\Delta t$, son constantes. (a) Si $d$ se hace tres veces más grande, ¿la ecuación predice que $\Delta t$ se hará más grande o más pequeña? ¿En qué factor? (b) ¿Qué patrón de proporcionalidad de $\Delta t$ a $d$ predice la ecuación? (c) Para mostrar esta proporcionalidad como una línea recta en una gráfica, ¿qué cantidades debe graficar en los ejes horizontal y vertical? (d) ¿Qué expresión representa la pendiente teórica de esta gráfica?
